\clearpage\thispagestyle{empty}
\pdfbookmark[0]{Sources for Vesta (24 December 1812)}{vesta1812dec24-sources}
\begin{center}\textbf{Sources for Vesta (24 December 1812)}\end{center}
\noindent This AI translation follows the report of 24 December 1812 in Gauss’s \href{https://gdz.sub.uni-goettingen.de/id/PPN236007467}{\textit{Werke}, vol.~VI} (Göttingen, 1874), pp.~361--364. Its first printing is headed Göttingen in \textit{Göttingische gelehrte Anzeigen}, no.~205, pp.~\href{https://api.digitale-sammlungen.de/iiif/image/v2/bsb10538682_01007/full/1600,/0/default.jpg}{2041}, \href{https://api.digitale-sammlungen.de/iiif/image/v2/bsb10538682_01008/full/1600,/0/default.jpg}{2042}, \href{https://api.digitale-sammlungen.de/iiif/image/v2/bsb10538682_01009/full/1600,/0/default.jpg}{2043}, \href{https://api.digitale-sammlungen.de/iiif/image/v2/bsb10538682_01010/full/1600,/0/default.jpg}{2044}, \href{https://api.digitale-sammlungen.de/iiif/image/v2/bsb10538682_01011/full/1600,/0/default.jpg}{2045}, here linked to the Bavarian State Library’s digitized copy. The report includes the closing account of Nicolai’s work on the second comet of 1811.
\par\smallskip\noindent Both editions leave the day unstated in the line giving the October opposition. In Enke’s elements on p.~2044, the first printing ends the longitude of perihelion with $31{,}1$ and includes degree and minute marks and a seconds mark in the sine expression for the eccentricity. The collected text has $249^\circ\;38^{\prime}\;31.4$ for the longitude and gives the sine argument without those marks. Translation and typesetting produced by AI.
\clearpage
